\documentclass[11pt]{amsart}

\usepackage{amsmath, amssymb, amsthm}
\usepackage{mathtools}
\usepackage{thmtools}
\usepackage{enumitem}
\usepackage{tikz-cd}
\usepackage{hyperref}
\usepackage[nameinlink]{cleveref}

% % Theorem environments
\declaretheorem[numberwithin=section]{theorem}
\declaretheorem[style=plain, sibling=theorem]{lemma, proposition, corollary}
\declaretheorem[style=definition]{definition, example, condition}
\declaretheorem[style=remark]{remark, note, observation, todo}


\title{April 15, 2026}


\begin{document}
\maketitle

After daily exercise and recording playing drum, we checked the feedback file for the yesterday diary.
The main feedback was as following:
\begin{itemize}
    \item Article usage,
    \item Tense consistency,
    \item To use simpler verb phrases,
    \item To split bibliographic context and mathematical idea,
    \item To make sentence surrounding sentences clear,
    \item To use ``my understanding is'' for tentative mathematical interpretations.
\end{itemize}

\par\vspace{1em}
After having lunch, we tested the skill \texttt{arXiv-review-routine}.
As there was a DNS error again, we decided to fix this problem.
The cause of DNS failure was that a sandbox executing Codex's automation was not allowed for networ access.
What's dissapointing was that we spent too much time on debugging.
We wondered for figuring out what is real problem, and how to fix it.
Eventhough I only gave an order to Codex, there was time consumption on waiting for the report of Codex.

\par\vspace{1em}
After having break by chatting with other PostDoc in our university, we continued to reading McPhail-Snyder's paper from the second section.
This section focuses on the viewpoint that the classical Chern--Simons invariant can be seen as a line bundle.
The author persuade readers that, the quantum invariant constructed in his paper is a vector bundle, and that it seems like the quantization of the CS invariant.
Here we record the current understanding.

\par\vspace{1em}

Let $\mathcal{M} := \mathbb{C}^2$, the universal cover of $\mathcal{X} = (\mathbb{C}^{\times})^2$ with the projection map which is the exponential map.
Regarding the trivial bundle $\mathcal{M} \times \mathbb{C}$, one might regard a section of this bundle as a functional on $\mathcal{M}$.
Moreover, if we can extend the action of $\mathbb{Z}^2 \cong \pi_1(\mathcal{X})$ on $\mathcal{M}$ to this bundle, then we have a quotient bundle $(\mathcal{M} \times \mathbb{C})/\mathbb{Z}^2$ over $\mathcal{X}$.
In other words, we have the following commutative diagram.
\[
    \begin{tikzcd}
    M \times \mathbb{C} \arrow[r, two heads] \arrow[d, two heads]
    & (M \times \mathbb{C})/\mathbb{Z}^2 \arrow[d] \\
    M \arrow[r, two heads]
    & X
    \end{tikzcd}
\]
and we regard \( M \times \mathbb{C} \) as a pull-back bundle of \( (M \times \mathbb{C})/\mathbb{Z}^2 \).

For any (holomorphic) section $f$ of \( (M \times \mathbb{C})/\mathbb{Z}^2 \) then we have a (holomorphic) pull-back section of $M \times \mathbb{C}$, equivalent to a functional on $M$ with the restriction induced by the action of $\mathbb{Z}^2$.

In particular, the action considered in the reference is that
\[
    (a,b) \cdot (x, y, z) = (x+a, y+b, e^{2\pi i(bx - ay)} z),
\]
hence the section $f$ satisfies that 
\[
    f(x+a, y+b) = e^{2\pi i(bx - ay)} f(x,y),
\]
which is called the \emph{quasi-periodicity relation}.



\bibliographystyle{amsalpha}
% \nocite{*}
\bibliography{bibliography}
\end{document}
